% Options for packages loaded elsewhere
\PassOptionsToPackage{unicode}{hyperref}
\PassOptionsToPackage{hyphens}{url}
\PassOptionsToPackage{dvipsnames,svgnames,x11names}{xcolor}
%
\documentclass[
  10pt,
  a4paper,
]{scrreprt}

\usepackage{amsmath,amssymb}
\usepackage{iftex}
\ifPDFTeX
  \usepackage[T1]{fontenc}
  \usepackage[utf8]{inputenc}
  \usepackage{textcomp} % provide euro and other symbols
\else % if luatex or xetex
  \usepackage{unicode-math}
  \defaultfontfeatures{Scale=MatchLowercase}
  \defaultfontfeatures[\rmfamily]{Ligatures=TeX,Scale=1}
\fi
\usepackage{lmodern}
\ifPDFTeX\else  
    % xetex/luatex font selection
\fi
% Use upquote if available, for straight quotes in verbatim environments
\IfFileExists{upquote.sty}{\usepackage{upquote}}{}
\IfFileExists{microtype.sty}{% use microtype if available
  \usepackage[]{microtype}
  \UseMicrotypeSet[protrusion]{basicmath} % disable protrusion for tt fonts
}{}
\usepackage{xcolor}
\usepackage[inner=2cm,outer=2cm,top=2cm,bottom=2cm,headsep=22pt,headheight=11pt,footskip=33pt,ignorehead,ignorefoot,heightrounded]{geometry}
\setlength{\emergencystretch}{3em} % prevent overfull lines
\setcounter{secnumdepth}{1}
% Make \paragraph and \subparagraph free-standing
\ifx\paragraph\undefined\else
  \let\oldparagraph\paragraph
  \renewcommand{\paragraph}[1]{\oldparagraph{#1}\mbox{}}
\fi
\ifx\subparagraph\undefined\else
  \let\oldsubparagraph\subparagraph
  \renewcommand{\subparagraph}[1]{\oldsubparagraph{#1}\mbox{}}
\fi


\providecommand{\tightlist}{%
  \setlength{\itemsep}{0pt}\setlength{\parskip}{0pt}}\usepackage{longtable,booktabs,array}
\usepackage{calc} % for calculating minipage widths
% Correct order of tables after \paragraph or \subparagraph
\usepackage{etoolbox}
\makeatletter
\patchcmd\longtable{\par}{\if@noskipsec\mbox{}\fi\par}{}{}
\makeatother
% Allow footnotes in longtable head/foot
\IfFileExists{footnotehyper.sty}{\usepackage{footnotehyper}}{\usepackage{footnote}}
\makesavenoteenv{longtable}
\usepackage{graphicx}
\makeatletter
\def\maxwidth{\ifdim\Gin@nat@width>\linewidth\linewidth\else\Gin@nat@width\fi}
\def\maxheight{\ifdim\Gin@nat@height>\textheight\textheight\else\Gin@nat@height\fi}
\makeatother
% Scale images if necessary, so that they will not overflow the page
% margins by default, and it is still possible to overwrite the defaults
% using explicit options in \includegraphics[width, height, ...]{}
\setkeys{Gin}{width=\maxwidth,height=\maxheight,keepaspectratio}
% Set default figure placement to htbp
\makeatletter
\def\fps@figure{htbp}
\makeatother

\addtokomafont{disposition}{\rmfamily}
\usepackage{ragged2e}
\usepackage{blindtext}\usepackage{amsthm}
\usepackage{amsthm}
\newtheorem{thm}{Theorem}[subsection]
\renewcommand{\thethm}{\arabic{subsection}.\arabic{thm}}
\newcommand{\GEV}{\text{GEV}}
\newcommand{\GP}{\text{GP}}
\newcommand{\IGP}{\text{IGP}}
\makeatletter
\makeatother
\makeatletter
\makeatother
\makeatletter
\@ifpackageloaded{caption}{}{\usepackage{caption}}
\AtBeginDocument{%
\ifdefined\contentsname
  \renewcommand*\contentsname{Table of contents}
\else
  \newcommand\contentsname{Table of contents}
\fi
\ifdefined\listfigurename
  \renewcommand*\listfigurename{List of Figures}
\else
  \newcommand\listfigurename{List of Figures}
\fi
\ifdefined\listtablename
  \renewcommand*\listtablename{List of Tables}
\else
  \newcommand\listtablename{List of Tables}
\fi
\ifdefined\figurename
  \renewcommand*\figurename{Figure}
\else
  \newcommand\figurename{Figure}
\fi
\ifdefined\tablename
  \renewcommand*\tablename{Table}
\else
  \newcommand\tablename{Table}
\fi
}
\@ifpackageloaded{float}{}{\usepackage{float}}
\floatstyle{ruled}
\@ifundefined{c@chapter}{\newfloat{codelisting}{h}{lop}}{\newfloat{codelisting}{h}{lop}[chapter]}
\floatname{codelisting}{Listing}
\newcommand*\listoflistings{\listof{codelisting}{List of Listings}}
\usepackage{amsthm}
\theoremstyle{definition}
\newtheorem{definition}{Definition}[section]
\theoremstyle{plain}
\newtheorem{theorem}{Theorem}[section]
\theoremstyle{remark}
\AtBeginDocument{\renewcommand*{\proofname}{Proof}}
\newtheorem*{remark}{Remark}
\newtheorem*{solution}{Solution}
\makeatother
\makeatletter
\@ifpackageloaded{caption}{}{\usepackage{caption}}
\@ifpackageloaded{subcaption}{}{\usepackage{subcaption}}
\makeatother
\makeatletter
\@ifpackageloaded{tcolorbox}{}{\usepackage[skins,breakable]{tcolorbox}}
\makeatother
\makeatletter
\@ifundefined{shadecolor}{\definecolor{shadecolor}{rgb}{.97, .97, .97}}
\makeatother
\makeatletter
\makeatother
\makeatletter
\makeatother
\ifLuaTeX
  \usepackage{selnolig}  % disable illegal ligatures
\fi
\IfFileExists{bookmark.sty}{\usepackage{bookmark}}{\usepackage{hyperref}}
\IfFileExists{xurl.sty}{\usepackage{xurl}}{} % add URL line breaks if available
\urlstyle{same} % disable monospaced font for URLs
\hypersetup{
  pdftitle={Annual Progress Review},
  pdfauthor={Thomas William Boughen},
  colorlinks=true,
  linkcolor={blue},
  filecolor={Maroon},
  citecolor={Blue},
  urlcolor={Blue},
  pdfcreator={LaTeX via pandoc}}

\title{Annual Progress Review}
\author{Thomas William Boughen}
\date{}

\begin{document}
\cleardoublepage
\thispagestyle{empty}
{\centering
\hbox{}\vskip 0cm plus 1fill
{\Huge\bfseries Annual Progress Review \par}
\vspace{12ex}
{\Large\bfseries Thomas William Boughen \par}
\vspace{3ex}
\vskip 0cm plus 2fill
%{\bfseries\large Doctor of Philosophy \par}
\vspace{3ex}
{\bfseries\large  \par}
\vspace{12ex}
{\includegraphics[width=0.1\linewidth]{"imgs/University_of_Newcastle_Coat_of_Arms.png"}\par}
%
%
{\bfseries\large Newcastle University \par}
\vspace{3ex}
%
{\bfseries\large School of Mathematics, Statistics and Physics \par}
%
\vspace{12ex}
%{\small Submitted in total fulfilment of the requirements
%of the degree of Doctor of Philosophy \par}
%}
\justifying
\noindent\ifdefined\Shaded\renewenvironment{Shaded}{\begin{tcolorbox}[sharp corners, frame hidden, borderline west={3pt}{0pt}{shadecolor}, enhanced, boxrule=0pt, interior hidden, breakable]}{\end{tcolorbox}}\fi

\hypertarget{intuition}{%
\chapter*{Intuition}\label{intuition}}
\addcontentsline{toc}{chapter}{Intuition}

\hypertarget{extremes}{%
\chapter{Extremes}\label{extremes}}

Since the aim is to gain understanding about the behaviour of the degree
distribution of networks at the right tail, it seems natural to look to
using methods from extreme value theory. However, networks by their
nature are discrete and so it may not be best to be using methods that
are usually used in relation to continuous random variables. For this
reason, this section starts with a review of what theory exists for
modelling the extreme values of continuous random variables before
moving to details what can be used when instead considering discrete
random variables as is the case for the degree distributions of random
networks.

\hypertarget{continuous-extremes}{%
\section{Continuous Extremes}\label{continuous-extremes}}

Studying the properties of the right tail of the distribution of a
continuous random variable, means that the focus is on the largest
values that the random variable can take. So, a natural place to start
is to consider the distribution of the block maxima of such a random
variable. That is, for a set of iid random variables
\(\{X_1,\ldots,X_n\}\) with common cumulative density function (cdf)
\(F\) what is the distribution of \(M_n = \max\{X_1,\ldots,X_n\}\)? This
question is answered by the Fisher--Tippett--Gnedenko theorem
\href{or\%20more\%20simply\%20the\%20extreme\%20value\%20theorem}{REF}.

\begin{theorem}[Extreme Value
Theorem]\protect\hypertarget{thm-evt}{}\label{thm-evt}

Let \(X_1,\ldots,X_n\) be a sample of iid random variables with common
cdf \(F\) with block maxima \(M_n = \max\{X_1,\ldots,X_n\}\) and suppose
that there exists \(a_n>0, b_n\in\mathbb R\) such that
\(\lim_{n\rightarrow\infty}\Pr(\frac{1}{a_n}[M_n-b_n]) = G(x)\), then
\(F\) is said to be in the domain of attraction of \(G\), denoted
\(F\in\mathcal D(G)\) ,and \(G\) is of one of three types:

\begin{itemize}
\tightlist
\item
  Gumbel: \(\Lambda(x) = \exp\{-\exp(-x)\},\quad x \in \mathbb R\)
\item
  Fréchet:
  \(\Phi_\alpha(x) = \exp\{-x^{-\alpha}\},\quad x\ge 0,\alpha>0\)
\item
  Weibull: \(\Psi_\alpha(x) = \exp\{-x^{-a}\},\quad x<0,\alpha>0\)
\end{itemize}

\end{theorem}

While this is a useful result, it may prove difficult to find the
sequences \(a_n,b_n\) in practice, so a simpler method to establish what
domain of attraction a distribution belongs to would be nice. Luckily,
this can be done through the concept of regular variation and is what
will be used to define the domains of attraction and tail-heaviness
through the rest of this report.

\begin{definition}[Domains of
Attraction]\protect\hypertarget{def-doa}{}\label{def-doa}

The distribution \(F\) belongs to the Fréchet domain of attraction
\(\mathcal D(\Phi_\alpha)\) if and only if its complement (the survival
function) \(\bar F\) is regularly varying with index \(-\alpha\) i.e.:
\[
\bar F(x) = x^{-\alpha}L(x),\qquad \text{for } L \text{ slowly varying}
\] A similar condition applies to the Weibull domain of attraction
\(\mathcal D(\Psi_\alpha)\) in that a distribution \(F\) belongs to the
Weibull domain of attraction if and only if: \[
\bar F(x_F-x^{-1}) = x^{-\alpha}L(x),\qquad \text{for } L \text{ slowly varying}
\] where \(x_F\) is the finite right endpoint of the support of \(F\).

The condition for the Gumbel domain of attraction is not as simple, a
distribution \(F\) belongs to the Gumbel domain if and only if there
exists a positive function \(a: \mathbb R \rightarrow \mathbb R^+\) and
a \(t\in \mathbb R\) such that: \[
\lim_{x\rightarrow x_F} \frac{\bar F(x+ta(x))}{\bar F(x)} = e^{-t}
\]

\end{definition}

Throughout this report the term ``heavy tailed'' distribution will be
used to describe any distribution in the Fréchet domain of attraction,
although some of the literature refers to ``heavy tailed'' distributions
as being the distributions that decay slower than the exponential.

At this point it will also be useful to introduce the concept of
distributions that have super-heavy tails.

\begin{definition}[Super Heavy
Tails]\protect\hypertarget{def-sht}{}\label{def-sht}

{[}FIND SUPER HEAVY TAILS DEFINITION{]}

Additionally, if the survival function \(\bar F\) is slowly varying
itself then \(F\) has super heavy tails i.e. \[
\lim_{x\rightarrow\infty}\frac{\bar F(tx)}{\bar{F}(x)} = 1, \forall t\in\mathbb R^+ \implies \text{super heavy tails} 
\]

\end{definition}

It is possible to gather the main three types of extremal distributions
into what is called the Generalised Extreme Value (GEV) distribution
{[}REF{]}.

\begin{definition}[Generalised Extreme Value
Distribution]\protect\hypertarget{def-gev}{}\label{def-gev}

Denoted by \(\GEV(\mu,\sigma,\xi)\) the distribution is characterised by
three parameters \(\mu \in \mathbb R\) the location,
\(\sigma\in \mathbb R^+\) the scale, and the shape \(\xi\in \mathbb R\).
It has support on \(\{x\in \mathbb R:1+\xi(x-\mu)/\sigma > 0\}\) and has
cdf given by:

\[
G(x) = \begin{cases}\exp\left\{-\left(1+\frac{\xi(x-\mu)}{\sigma}\right)^{-1/\xi}\right\},&\xi\ne0\\
\exp\{-\exp(-\frac{x-\mu}{\sigma})\},&\xi=0
\end{cases}
\]

\end{definition}

The three types of extremal distribution are obtained from changing the
shape parameter \(\xi\), which corresponds to \(1/\alpha\) in the
definition of the domains of attraction. This change is generally to
made so that increasing \(\xi\) corresponds to increasing how heavy the
tails of the distribution are. So, \(\xi<0\) corresponds to the Weibull,
\(\xi>0\) the Fréchet, and \(\xi=0\) the Gumbel.

\hypertarget{discrete-extremes}{%
\section{Discrete Extremes}\label{discrete-extremes}}

\hypertarget{networks}{%
\chapter{Networks}\label{networks}}

\newpage{}

\hypertarget{references}{%
\chapter*{References}\label{references}}
\addcontentsline{toc}{chapter}{References}



\end{document}
